\documentclass[12pt]{article}
\usepackage{espd}
\docode{RU.МИФИ.52714-01 81 01-1}
% \makeapproval{Наименование}{Вид документа}{Шифр ЛУ}{Шифр}
\newcommand{\univ}{Национальный исследовательский ядерный университет «МИФИ»}
\newcommand{\course}{Дисциплина «Генеративный искусственный интеллект: модели и приложения»}
\newcommand{\progname}{Компьютерная игра «Dreich Day»}
\newcommand{\makeapproval}[3]{%
\thispagestyle{empty}
\begin{center}\singlespacing
\univ\\[2mm]\course
\end{center}
\vspace{10mm}
\noindent\begin{minipage}[t]{0.48\textwidth}\singlespacing
СОГЛАСОВАНО\\[2mm]Студент группы Б25-527\\[6mm]\rule{30mm}{0.4pt} А.\,И.~Лобанов\\[2mm]«\rule{8mm}{0.4pt}»\rule{30mm}{0.4pt} 2026 г.
\end{minipage}\hfill
\begin{minipage}[t]{0.48\textwidth}\singlespacing\raggedleft
УТВЕРЖДАЮ\\[2mm]Преподаватель дисциплины\\[6mm]\rule{30mm}{0.4pt} А.\,А.~Кузьмин\\[2mm]«\rule{8mm}{0.4pt}»\rule{30mm}{0.4pt} 2026 г.
\end{minipage}
\vfill
\begin{center}
{\bfseries\MakeUppercase{\progname}}\\[6mm]
{\bfseries\MakeUppercase{#1}}\\[4mm]
ЛИСТ УТВЕРЖДЕНИЯ\\[4mm]
#2
\end{center}
\vfill
\begin{flushright}\singlespacing
Исполнитель:\\студент группы Б25-527\\Лобанов Алексей Игоревич
\end{flushright}
\vspace{10mm}
\begin{center}Москва 2026\end{center}
\clearpage
\thispagestyle{empty}
\noindent УТВЕРЖДЁН\\#2\\
\vfill
\begin{center}
{\bfseries\MakeUppercase{\progname}}\\[6mm]
{\bfseries\MakeUppercase{#1}}\\[4mm]
#3\\[4mm]
Листов \pageref{LastPage}
\end{center}
\vfill
\begin{center}Москва 2026\end{center}
\clearpage
\setcounter{page}{2}}

\begin{document}
\makeapproval{Пояснительная записка}{RU.МИФИ.52714-01 81 01-1-ЛУ}{RU.МИФИ.52714-01 81 01-1}

\section*{\centering АННОТАЦИЯ}
В пояснительной записке описана компьютерная игра «Dreich Day», разработанная в рамках большого домашнего задания по дисциплине «Генеративный искусственный интеллект: модели и приложения». Приведены назначение и область применения программы, постановка задачи, описание алгоритма работы, организации данных, выбора технических и программных средств, а также описан процесс разработки с применением генеративного ИИ на всех этапах: от генерации идеи до написания кода (вайбкодинг на платформе WebSim) и создания графики.

Документ оформлен в соответствии с ГОСТ 19.404-79 «ЕСПД. Пояснительная записка. Требования к содержанию и оформлению».
\clearpage
\tableofcontents
\clearpage

\section{Введение}
\subsection{Наименование программы}
Компьютерная игра «Dreich Day».
\subsection{Документы, на основании которых ведётся разработка}
Задание на большое домашнее задание по дисциплине «Генеративный искусственный интеллект: модели и приложения» (преподаватель А.\,А.~Кузьмин); техническое задание RU.МИФИ.52714-01 ТЗ 01-1; спецификация требований RU.МИФИ.52714-01 СТ 01-1. Работа также представляется для зачёта в качестве проектной практики.

\section{Назначение и область применения}
Игра предназначена для развлечения и демонстрирует применение генеративного ИИ в полном цикле разработки программного продукта. Игрок в роли лондонского синоптика Оливера Оверкаста застревает во временной петле в шотландской деревне Глен-Дрих в День святой Бригиты. В отличие от фильма «День сурка», здесь каждый повтор дня складывается по-новому: меняются погода, события и настроение жителей. Задача игрока — понять законы, по которым живёт деревня, и подготовить её команду по шинти к финальному матчу против деревни Силлерберн, чтобы спасти любимый паб от сноса.

\section{Технические характеристики}
\subsection{Постановка задачи}
Требуется разработать браузерную игру, в которой: (1) день повторяется, но его условия генерируются случайно; (2) поведение жителей определяется скрытыми детерминированными правилами; (3) у игрока есть рычаг влияния — прогноз погоды в утренней газете; (4) прогресс внутри петли выражается в заполнении дневника правил и выполнении заданий; (5) кульминацией служит матч, выиграть который можно только при «идеальном дне».

Ключевой принцип: \textit{рандом выбирает условия дня, а постоянные правила превращают их в поведение}. Благодаря этому дни различаются, но их можно предсказать, если знать правила.

\subsection{Описание алгоритма и функционирования программы}
\subsubsection{Игровой цикл}
Основной цикл выполняется через \texttt{requestAnimationFrame}: на каждом кадре обновляется игровое время (1 игровой час = 50 с), положение игрока и жителей, проверяются правила, задания и события, затем выполняется отрисовка с сортировкой объектов по глубине.
\subsubsection{Генерация условий дня}
В 6:00 функция генерации выбирает погоду утром и после полудня (после полудня с вероятностью 50\,\% совпадает с утренней), знак гадюки на обряде (совпадает с погодой после полудня с вероятностью 70\,\%), событие дня и стартовое настроение жителей (40–70).
\subsubsection{Движок правил}
Правило — объект данных с идентификатором, жителем, условием \texttt{when(state)} и эффектом (место, интервал времени, изменение настроения). Правила переопределяют базовое расписание жителя. Жители перемещаются по общей сетке проходимости с поиском пути в ширину (BFS), входя в здания только через двери.
\subsubsection{Прогноз, дневник и задания}
Прогноз выбирается один раз за петлю и проверяется в 12:00. Наблюдение фиксируется, когда игрок находится рядом с жителем при срабатывании правила; после двух наблюдений в разных петлях правило открывается. Задание дня выбирается из пула с фильтрацией по выполнимости при текущих скрытых условиях.

Отношения между персонажами описаны данными: конфликт — пара жителей, эффект на настроение, пока конфликт не разрешён, и условия разрешения; история — цепочка шагов, которые сохраняются между петлями и отображаются в дневнике знаками «???» до первого прохождения. Реализованы истории «Старые друзья» (Дуги и Ангус) и «Подкупленный судья» (разоблачение Руперта Скуннерсби через экстренный выпуск газеты) и конфликты «Пикша и булки», «Псалмы и волынка». Разрешённые конфликты и истории дают новые совместные правила и постоянные флаги, которые учитываются в финальном матче. Финальный матч (итерация 4) начинается в 16:00, если Оливер пришёл на поле. На него приходят только жители с настроением не ниже 40. Сила команды равна сумме настроений вышедших игроков с бонусом за разрешённые конфликты, штрафом за неразрешённые, поправкой на погоду и прогноз и штрафом за подкупленного судью; отношение к силе Силлерберна задаёт скорость, точность и отбор мяча у игроков ИИ. Матч — мини-игра 6 на 6 длительностью 75 секунд, все числа баланса вынесены в один объект конфигурации. Итерация 5 добавила концовку (победа засчитывается, только если решающий гол забил Оливер), титры, главное меню, подсказки первой петли и простые звуки, синтезируемые через Web Audio без внешних файлов.
\subsubsection{Сброс петли}
В 18:00 или по решению игрока показывается сводка дня; дневник, счётчик петель и прогресс заданий сохраняются в localStorage, остальное состояние сбрасывается.

\subsection{Организация входных и выходных данных}
Входные данные — клавиатурный ввод (WASD/стрелки, E, J, Esc; используется \texttt{event.code}, поэтому управление работает при любой раскладке) и выбор прогноза мышью. Выходные данные — изображение на HTML5 Canvas (320×240 логических пикселей, масштаб ×3), всплывающие уведомления, панель задания, дневник. Сохранение — объект JSON в localStorage.

\subsection{Выбор состава технических и программных средств}
Платформа WebSim выбрана по условию задания: она генерирует код по описанию на естественном языке (вайбкодинг). Игра реализована на HTML5, CSS и JavaScript без внешних библиотек в трёх файлах (\texttt{index.html}, \texttt{style.css}, \texttt{game.js}). Графика выполнена в пиксельном стиле; спрайты персонажей, зданий, интерьеров, тайлы земли и горное окружение сгенерированы с помощью ИИ-ассистента Claude программно (попиксельно в заданной палитре) и подключены как ресурсы проекта.

\subsection{Применение генеративного ИИ в разработке}
\subsubsection{Этап 1. Мозговой штурм}
Идея получена в диалоге с ИИ: от исходной ассоциации с фильмом «День сурка» через ряд итераций (перенос в Шотландию, реальный фольклор Là Fhèill Brìghde о гадюке, предсказывающей погоду, конфликт деревни с лондонским дельцом, шинти как кульминация) к итоговой концепции. Персонажам даны говорящие шотландские имена (Overcast, Scunnersby, Bannock, Haddock и др.).
\subsubsection{Этап 2. Техническое задание}
ТЗ составлено в режиме мета-фасилитации: ИИ поочерёдно выступал заказчиком-геймером, архитектором и бизнес-аналитиком, исполнитель принимал решения как владелец продукта. Результат — ТЗ, план из пяти итераций и оптимизированный промпт для WebSim.
\subsubsection{Этап 3. Вайбкодинг на WebSim}
Разработка велась итерациями: крупный промпт с новой функциональностью, затем плейтест и короткие точечные запросы на исправление с явным ограничением «не менять остальное». Журнал запросов приведён в приложении.

\subsubsection{Этап 4. Документирование}
Комплект документов (техническое задание, спецификация требований, пояснительная записка) подготовлен в \LaTeX{} по ГОСТ 19.101-77 и ГОСТ 19.106-78 с помощью генеративного ИИ: модель составляла текст по материалам ТЗ, кода и журнала, исполнитель проверял и утверждал. Шаблон оформления (титульная часть, лист утверждения, колонтитул с обозначением, лист регистрации изменений) реализован отдельным стилевым файлом и общий для всех документов.

\subsubsection{Этап 5. Репозиторий}
Исходный код игры, все спрайты, \LaTeX-исходники и PDF документов, промпты всех итераций и скрипты генерации графики опубликованы в публичном репозитории (GitHub, зеркало на GitVerse).

\section{Особенности выполнения работы}
\begin{enumerate}
\item \textbf{Мета-фасилитация и метапромптинг.} Техническое задание получено в диалоге, где модель последовательно играла роли agile-команды (заказчик-геймер, архитектор, бизнес-аналитик); промпты для WebSim составлялись и оптимизировались генеративным ИИ: роль и цель в первой строке, точные числа, правила как данные, явный список того, что делать не нужно, ограничение «не менять остальное» в точечных запросах.
\item \textbf{Сложные механики.} Временная петля с сохранением знаний игрока, 26 детерминированных скрытых правил 11 жителей, задания с фильтром выполнимости, конфликты и истории, сохраняющиеся между петлями, финальный матч с формулой силы команды и ИИ игроков.
\item \textbf{Нестандартное применение генеративного ИИ.} Вся пиксельная графика (16 персонажей, 9 построек, нора, лодка, овцы, тайлсет, поле для шинти) нарисована генеративным ИИ программно — Python-скриптами на Pillow, без генераторов изображений; плейтест каждой версии выполнялся ИИ-агентом автоматически в браузере, найденные ошибки формулировались в точечные промпты.
\item \textbf{Сроки.} Идея и ТЗ — 5 октября 2026 г., пять итераций игры, документация и репозиторий — к 7 ноября 2026 г.
\end{enumerate}

\section{Ожидаемые технико-экономические показатели}
Разработка выполнена одним исполнителем на бесплатных ежедневных кредитах WebSim. Стоимость одного крупного запроса составила 75–315 кредитов, точечного исправления — 20–100 кредитов; всего на 18 запросов журнала израсходовано около 2100 кредитов.

\section{Источники, использованные при разработке}
\begin{enumerate}
\item ГОСТ 19.101-77 ЕСПД. Виды программ и программных документов.
\item ГОСТ 19.106-78 ЕСПД. Требования к программным документам, выполненным печатным способом.
\item ГОСТ 19.201-78 ЕСПД. Техническое задание. Требования к содержанию и оформлению.
\item ГОСТ 19.404-79 ЕСПД. Пояснительная записка. Требования к содержанию и оформлению.
\item Carmina Gadelica, vol. 1. — Национальная библиотека Шотландии (фольклор Là Fhèill Brìghde).
\item Документация платформы WebSim — websim.com.
\end{enumerate}

\clearpage
\section*{\centering ПРИЛОЖЕНИЕ 1. ЖУРНАЛ РАЗРАБОТКИ НА WEBSIM}
\addcontentsline{toc}{section}{Приложение 1. Журнал разработки на WebSim}
{\small\setlength{\tabcolsep}{3pt}\renewcommand{\arraystretch}{1.15}
\begin{longtable}{|p{6mm}|p{17mm}|p{62mm}|p{58mm}|p{11mm}|}\hline
№ & Дата & Запрос к агенту WebSim (суть) & Результат & Кре\-диты\\\hline\endhead
\multicolumn{5}{|l|}{\textbf{Итерация 1 — «Петля»}}\\\hline
1 & 05.10.2026 & Полный промпт итерации 1 из ТЗ (с блоком арт-дирекшена и палитрой) & Игра в трёх файлах: карта 40×30, погода, жители с правилами, прогноз, обряд, дневник, сводка дня; игрок не мог ходить & 105\\\hline
2 & 05.10.2026 & Исправить движение: спавн, \texttt{event.code}, коллизии по осям & Ходьба работает при любой раскладке & 25\\\hline
3 & 05.10.2026 & Плейтест-фидбек (10 пунктов): поиск пути через двери, обряд без телепортации, стук в дверь, подсказки E, уведомления о наблюдениях, подсказки в репликах и др. & Все пункты реализованы & 140\\\hline
4 & 05.10.2026 & Прогноз один раз за петлю, жители не стоят на одной клетке, счётчик дневника & Итерация 1 принята и опубликована & 40\\\hline
\multicolumn{5}{|l|}{\textbf{Итерация 2 — «Цели»}}\\\hline
5 & 06.10.2026 & Задания дня: пул из 5 заданий с фильтром выполнимости, объявление в газете, награда — подсказка к правилу & Утренняя газета, панель задания, раздел «Задания» в дневнике & 105\\\hline
6 & 06.10.2026 & Добавить подсказку и цель задания & Подсказка в газете, цель в панели & 20\\\hline
7 & 06.10.2026 & Замена процедурной графики на спрайты персонажей и зданий & Анимация по направлениям, сортировка по глубине, прозрачность здания при входе & 85\\\hline
8 & 06.10.2026 & Горное окружение карты & Горы, лес, скалы, озеро по краям; камера ограничена миром & 75\\\hline
9 & 06.11.2026 & Тайлсет земли (трава, дороги, брусчатка, вода, берег) и исправление окружения & Новые тайлы, горы видны полностью, лес не растёт на воде & 45\\\hline
10 & 06.11.2026 & Новые жители (Дуги, Элспет, Иэн), визиты Ангуса, кирк, домик Дуги, поле для шинти, лодка Хэмиша, 3 новых задания & 6 жителей на обряде, 8 заданий; дороги слишком широкие, имена налезают & 180\\\hline
11 & 06.11.2026 & Дороги в одну клетку, мощёная площадь, одна табличка имени, исправление реплик & Карта аккуратная, ошибок в консоли нет & 60\\\hline
\multicolumn{5}{|l|}{\textbf{Итерация 2.3 — «Конфликты и истории»}}\\\hline
12 & 06.11.2026 & Конфликты как данные, первая история «Старые друзья», раздел «Истории» в дневнике, наблюдения только на месте & Сюжет из трёх шагов, сцена в пабе, новое совместное правило & 135\\\hline
\multicolumn{5}{|l|}{\textbf{Итерация 3 — «Вся деревня»}}\\\hline
13 & 06.11.2026 & Пять новых жителей (Рэб, Кенни, Малкольм, Рори, Бриди), кузница, сторожка, хижина Бриди, пастбище с овцами, 2 новых задания & Все 11 жителей на обряде, овцы бродят в загоне, ошибок нет; на обряде толпа закрывает нору & 295\\\hline
14 & 06.11.2026 & Антагонист Руперт (событие, история «Подкупленный судья»: подслушать, квитанция, экстренный выпуск), конфликты «Пикша и булки» и «Псалмы и волынка», правило Бриди B2, круг жителей вокруг норы & История проходится полностью, заголовок в газете, нора видна; ошибок нет & 315\\\hline
\end{longtable}}


\noindent Примечание. В период между запросами 8 и 9 в проект вносились дополнительные мелкие правки (версии до v17), в журнал они не заносились.

\changelist
\end{document}
